\documentclass[11pt]{article}
\usepackage[letterpaper,margin=1.04in]{geometry}
\usepackage{amsmath,amssymb,amsthm}
\usepackage[hidelinks]{hyperref}
\newtheorem{theorem}{Theorem}
\newtheorem{lemma}[theorem]{Lemma}
\newcommand{\TC}{\mathrm{TC}^{0}}
\newcommand{\FO}{\mathrm{FO}}
\newcommand{\Eq}{\operatorname{Eq}}
\newcommand{\Add}{\operatorname{Add}}
\newcommand{\Pow}{\operatorname{Pow}}
\newcommand{\Mul}{\operatorname{Mul}}
\newcommand{\Bit}{\operatorname{Bit}}
\newcommand{\Sh}{\operatorname{Shift}}
\newcommand{\Res}{\operatorname{Res}}
\newcommand{\dep}{C_E}
\title{Power-of-two counting captures uniform \(\TC\)}
\author{Mira}
\date{November 2026}

\begin{document}
\maketitle
\begin{abstract}
Let \(C_E x\,\varphi\) hold if the set defined by \(\varphi\) has
power-of-two cardinality. On finite linear orders, we define the ternary
relation \(ab=c\) by a single first-order formula with at most three nested
occurrences of \(C_E\), without enlarging the order. It follows that
\(\FO_{\le}(C_E)\) captures DLOGTIME-uniform \(\TC\) on finite ordered
words, answering a question of Hella, Kontinen, and Luosto. The key
operation shifts a binary integer by counting a union of disjoint
dyadic intervals. The largest in-domain weight is handled separately
so that the construction works up to the last element of the order.
\end{abstract}

\section{Statement}
Let \(E=\{2^i:i\geq 0\}\). The quantifier \(C_E x\,\varphi\)
tests the cardinality of the set of elements satisfying \(\varphi\),
with all other variables treated as parameters. The empty set fails
the test. We write \([n]=\{0,\ldots,n-1\}\), with its standard order.
The depth of a formula means its maximum nesting of \(C_E\) alone,
after all abbreviations have been expanded. Ordinary first-order
quantifiers do not add to this depth.

\begin{theorem}\label{thm:main}
There is a formula of depth at most three in \(\FO_{\le}(C_E)\)
that defines \(ab=c\) for all \(a,b,c\in[n]\) and every \(n\geq 1\).
In particular, on nonempty finite ordered words,
\(\FO_{\le}(C_E)\) captures DLOGTIME-uniform \(\TC\).
\end{theorem}

Hella, Kontinen, and Luosto ask this capture question in
\cite[Section~6]{hkl}. They give equicardinality for this quantifier
and a partial-multiplication construction for suitable counting
quantifiers. We use equicardinality with explicit depth bookkeeping,
then give a small-product relation using separated intervals and
congruences. The new step is the shift relation in Section~\ref{s:shift},
which removes the need to enlarge the order.

\section{Equicardinality at one counting layer}\label{s:eq}
Use \(\Eq(A,B)\) for a formula comparing the cardinalities of two
unary definable sets. We first exhibit a macro that adds at most one
layer to the depth of the two input predicates. Set expressions in
the proof are substitutions of unary formulas, never quantification
over arbitrary sets.

For \(n\geq4\), find the largest power-sized prefix and suffix,
each of size \(4P\). They are definable by testing prefixes and
suffixes with \(C_E\) and maximizing under inclusion.
They cover \([n]\), since \(4P\leq n<8P\).
An interval of size \(2^d\) splits into equal halves at the unique cut
with both sides of power-of-two size: two positive powers can sum
to a power only when equal. Split the prefix and suffix twice in
this way. The four prefix quarters and the four suffix quarters
after removing the prefix are disjoint pieces of size at most \(P\).
All cuts use ordinary endpoint witnesses and parallel \(C_E\) tests.
The first quarter and its immediate successor represent \(0,\ldots,P\).

For sets \(A,B\) of size at most \(P\), let \(W\) be the
complement of their union. Compare the truth values of the two power
tests on \(A\cup W'\) and \(B\cup W'\) for every nonempty initial
segment \(W'\) of \(W\). Equality of sizes gives agreement. If the
sizes differ, take \(|W'|=2P-\min(|A|,|B|)\), which is possible
since \(|W|\geq 2P\). One total is \(2P\), and the other lies
strictly between \(2P\) and \(4P\). This detects inequality.
Each \(W'\) is specified by the element of \(W\) at which it ends.

For arbitrary sets, compare the intersection with each of the eight
pieces to a prefix of size \(s\leq P\) by the small-set test.
To add these eight counts, use quotient and remainder in base \(P\),
with a remainder \(r<P\) and at most one carry at each of the eight
steps. Small addition \(i+j=k\leq P\) compares \([i,k)\) and
\([0,j)\) with the small-set test, requiring \(i\leq k\).
For a carry from \(r+s\) to \(r'\), require \(r+t=P\) and
\(t+r'=s\). The two lists of carries have the same number of ones
by a fixed Boolean condition; their final remainders must agree.
This works also for \(P=1\) and when a set is the whole domain.
For \(n\leq3\), enumerate the four possible cardinalities with
ordinary first-order formulas. This proves the claimed depth bound.

Exact truncated addition \(\Add(a,b,c)\) is \(a\leq c\) and
\(\Eq([a,c),[0,b))\), of depth one. Write \(P_2(p)\) for
\(C_E x\,(x<p)\); it says that \(p\) is a positive power of two.
The indexed power relation
\[
 \Pow(i,p)\ \equiv\ P_2(p)\ \land\
   \Eq(\{u:u<p,\ P_2(u)\},[0,i))
\]
has depth two and says \(p=2^i<n\). In particular \(p=1\)
corresponds to \(i=0\).

\section{A small product at depth two}\label{s:small}
We use an auxiliary relation \(D(x,y,z)\). It is sound for
\(z=xy\) on every order but initially only complete with sufficient
space. The zero-factor cases are immediate, so take \(x,y>0\).
Let \(p\) be the least in-domain power of two at least \(x\), if
there is one. The allowed starts are \(0\), the powers of two at
least \(p\), and sums of two distinct such powers. Select the starts
up to a cutoff which is itself a start. Compare this set to
\([0,y)\) by \(\Eq\), require that the cutoff plus \(x-1\) is still
in the domain, and compare the union of the length-\(x\) intervals
at the selected starts to \([0,z)\) by \(\Eq\).

All these starts are distinct multiples of \(p\), so the intervals
are disjoint. A successful test gives \(z=xy\). Membership in
the start or interval set uses only first-order operations,
\(P_2\), and truncated addition in parallel. The added layer
for \(\Eq\) makes the depth at most two. If \(t\) bit positions
starting at \(p\) fit and \(1+t(t+1)/2\geq y\), there are
at least \(y\) starts no larger than \(p(2^t-1)\). Thus the
test succeeds for a true product (with its arguments in the
domain) if
\[
 n\ \geq\ p(2^t-1)+x,\qquad 1+t(t+1)/2\geq y.
 \tag{1}\label{eq:space}
\]
The first inequality also ensures that the \(t\) positions used here
fit.

Define \(\Res(x,m,u)\), for \(m\geq2\), by \(u<m\) and
existence of \(v,w\) with \(D(v,m,w)\) and \(w+u=x\).
Whenever this succeeds, \(u\) is the genuine residue of \(x\).
For nonzero factors, define \(\Mul(a,b,c)\) by the following
congruence tests. Let \(w=\max(a,b,c)\), let \(e\) be the number
of powers of two at most \(w\), and put \(K=2e+2\).
The integer \(e\) is defined by equicardinality with \([0,e)\);
\(K\) uses truncated addition. For every \(2\leq m\leq K\),
find the residues of \(a,b,c\) using \(\Res\), multiply the
first two residues using \(D\), and require the residue of that
product to equal the residue of \(c\). Once again check zero
factors separately. Every arithmetic call is in ordinary
first-order scope, so \(\Mul\) has depth at most two.
Since \(w<2^e\), we have \(|ab-c|<2^{2e}\), and the tests force
\(ab=c\) by
\[
 \operatorname{lcm}(1,\ldots,K)\ \geq\ 2^{K-1}.
\]
The latter is the standard consequence of Farhi's identity
\cite[Theorem~1]{farhi}: \(\operatorname{lcm}(1,\ldots,K)\) is
\(K\) times the least common multiple of the binomial coefficients
in row \(K-1\); the largest coefficient is at least \(2^{K-1}/K\).

\begin{lemma}\label{lem:small}
Suppose \(n\geq65537\). Let \(2^h\) be the largest power of two in
\([n]\), let \(k=\lfloor h/4\rfloor\), and put \(B=2^k\).
For \(a,b,c<B^2\), the relation \(\Mul(a,b,c)\) holds exactly
when \(ab=c\).
\end{lemma}
\begin{proof}
Soundness was proved above. Here \(k\geq4\) and \(K\leq4k+2\).
In the true residue calculations and the product of two residues,
every \(D\) has second argument at most \(K\). Its first argument
is at most \(2B^2\): the original arguments are below \(B^2\), and
the residue product is bounded by \((4k+2)^2\leq2B^2\).
The dyadic ceiling of any positive width is therefore at most
\(2B^2\). With \(t=2k-2\), the number \(1+t(t+1)/2\) is at least
\(4k+2\), and the space in \eqref{eq:space} is bounded by
\[
 (2B^2)\,2^{2k-2}=2^{4k-1}<n.
\]
Consequently every required witness fits and all true tests succeed.
\end{proof}

\section{Counting a binary shift}\label{s:shift}
On the guarded range \(0\leq q<B^2\), define \(\Bit(q,j)\) by
the existence of \(p,t,r,y\) with
\[
 \Pow(j,p),\quad t>0\text{ odd},\quad r<p,\quad
 \Mul(t,p,y),\quad \Add(y,r,q).
\]
Oddness is expressible by truncated addition and successor. The
formula is sound by division with remainder. If bit \(j\) is set,
all witnesses, including \(p\), are at most \(q\); Lemma~\ref{lem:small}
then proves completeness. The depth is two.

To define \(\Sh(q,e,z)\) on the same input range, require that every
set bit \(j\) of \(q\) has a represented index \(e+j\) and a represented
target weight \(2^{e+j}\). Let \(P=2^h\) be the largest in-domain
power. Call a weight selected if it is the target of some set bit.
For each selected \(p<P\), use the interval \([p,2p)\), located
using \(\Add(p,p,2p)\). These intervals are disjoint and contained
in \([1,P)\). By \(\Eq\), require their union to have size exactly
\(s\) by comparing it with \([0,s)\). If \(P\) was selected, require
\(\Add(s,P,z)\); otherwise require \(s=z\).

The separate top term prevents clipping at the endpoint \(n\).
The size \(s<P\) always fits, including for \(q=0\). If \(q=0\)
the relation gives \(z=0\) for every shift exponent in the domain,
even if \(2^e\) is not in the domain. If \(q>0\) and the true output
is in the domain, each target is in the domain; its index is also
representable. Conversely, if a target is missing, or if the final
addition would overflow, the formula fails. Hence the relation
is exactly \(z=q2^e<n\) on the guarded range. Union membership
has depth two, because the power, bit, and addition subformulas
occur in parallel. The outer equicardinality macro brings the
depth to at most three.

\section{Full multiplication and capture}
For \(n\geq65537\), choose \(h,k,B\) as above. They are definable
using indexed powers, addition, and the remainder of \(h\) modulo
four, a finite choice. The cutoff \(B^2\) is also defined by
\(\Pow(2k,B^2)\). Since \(h\leq4k+3\) and \(k\geq4\), we have
\(B^5\geq2^{h+1}\geq n\). Thus five base-\(B\) digits suffice.

Existentially choose digits \(d_i,f_i<B\), \(0\leq i\leq4\).
Use \(\Sh\) and fixed chains of truncated additions to require
\[
 a=\sum_{i=0}^4 d_i\,2^{ik},\qquad b=\sum_{i=0}^4 f_i\,2^{ik}.
\]
For each pair \(i,j\), require \(g_{ij}<B^2\) and
\(\Mul(d_i,f_j,g_{ij})\). Use \(\Sh\) and addition again to demand
\[
 c=\sum_{i,j=0}^4 g_{ij}\,2^{(i+j)k}.
\]
The notation is a finite template: all indices are unrolled, and
the exponents, at most \(8k<n\), are defined by fixed additions.
Each witnessed equality is genuine. For a true product \(ab=c<n\),
each nonnegative shifted term and each partial sum fit below \(n\),
as do the terms and sums reconstructing the inputs. Thus the formula
is complete. Combining the finitely many calls in ordinary first-order
scope keeps depth at most three. The finitely many smaller orders
can be covered by a pure-order disjunction specifying their sizes
and the ordinal positions of each true triple.

By \cite[Section~2, especially equation~(5)]{hkl}, order,
multiplication, and H\"artig equicardinality capture DLOGTIME-uniform
\(\TC\) on nonempty ordered words. Both of the latter relations
are available in our logic. For the reverse inclusion, in that
arithmetic logic compare a nonempty set with the prefix ending
inclusively at \(u\). Its size is a power of two just when
\(u=2^i-1\) for some \(i\geq0\). This is the usual binary test
that \(u\) has no zero bit below a one bit, with \(u=0\) permitted.
Here BIT is definable from the built-in arithmetic
\cite[Chapter~6]{libkin}. The inclusive prefix also handles a
power-sized full domain, without representing \(n\) as an element.
This proves both inclusions and Theorem~\ref{thm:main}.

\paragraph{The remaining depth question.}
Can exact multiplication be defined with at most two nested
\(C_E\) quantifiers? Depth one cannot suffice even with built-in
truncated addition. Indeed, relativize each such formula to
\([0,n)\) in the natural numbers. Each counting child is
Presburger-definable; its unary cardinality is also
Presburger-definable, by \cite{schweikardt}. Testing that cardinality
for membership in \(E\) translates the whole formula into
Presburger arithmetic with a power-of-two predicate, a decidable
theory \cite{benedikt}. A depth-one definition of exact multiplication
would define multiplication on the natural numbers by existentially
choosing a containing \(n\), contradicting undecidability.
The distinction between depth two and three remains open here.

\paragraph{Acknowledgment.}
OpenAI models assisted in developing the proof, drafting the
exposition, and checking the argument, under Mira's guidance.
The proof was separately checked by a nonauthor automated reviewer;
no formal verification or external referee review is claimed.
A companion \href{https://github.com/Mira-Cult/math/tree/main/notes/logic/power-of-two-counting}{research writeup}
is available in the public mathematics repository.

\begin{thebibliography}{9}
\bibitem{hkl}
L. Hella, J. Kontinen, and K. Luosto,
\emph{Regular representations of uniform \(\TC\)},
ACM Transactions on Computational Logic 26 (2025), Article 21.
\href{https://doi.org/10.1145/3750044}{doi:10.1145/3750044}.

\bibitem{farhi}
B. Farhi,
\emph{An identity involving the least common multiple of binomial
coefficients and its application},
American Mathematical Monthly 116 (2009), 836--839.
\href{https://arxiv.org/abs/0906.2295}{arXiv:0906.2295}.

\bibitem{libkin}
L. Libkin, \emph{Elements of Finite Model Theory}, Springer, 2004.
\href{https://www.cs.toronto.edu/~libkin/fmt/fmt.pdf}{Online edition}.

\bibitem{schweikardt}
N. Schweikardt,
\emph{Arithmetic, first-order logic, and counting quantifiers},
ACM Transactions on Computational Logic 6 (2005), 634--671.
\href{https://arxiv.org/abs/cs/0211022}{arXiv:cs/0211022}.

\bibitem{benedikt}
M. Benedikt, C.-H. Chistikov, and A. Mansutti,
\emph{The complexity of Presburger arithmetic with power or powers},
ICALP 2023, Article 112.
\href{https://doi.org/10.4230/LIPIcs.ICALP.2023.112}{doi:10.4230/LIPIcs.ICALP.2023.112}.
\end{thebibliography}
\end{document}
